\chapter{Conception du modèle décisionnel}
\label{ch:modele}

\section{Introduction}
Conformément à la deuxième étape de la méthodologie GIMSI (conception), ce chapitre présente la modélisation multidimensionnelle retenue pour l'entrepôt de données. Le choix s'est porté sur un \textbf{schéma en constellation}, c'est-à-dire un modèle comportant plusieurs tables de faits partageant un ensemble commun de dimensions, ce qui correspond naturellement aux trois processus métier identifiés~: le comptage de trafic, les déplacements individuels et les caractéristiques des ménages.

\section{Choix du schéma en constellation}
Contrairement à un schéma en étoile classique (une seule table de faits), le schéma en constellation permet de modéliser plusieurs processus métier distincts tout en mutualisant les dimensions communes (temps, géographie). Ce choix est justifié par la nature du projet~: les données de comptage de trafic (issues de capteurs) et les données d'enquête ménages-déplacements ont des granularités et des cycles de mise à jour différents, mais partagent des axes d'analyse communs (site, zone géographique, temps).

Le modèle retenu comporte~:
\begin{itemize}
\item \textbf{3 tables de faits}~: \texttt{Fait\_Comptage}, \texttt{Fait\_Deplacement}, \texttt{Fait\_IndividuMenage}~;
\item \textbf{7 dimensions}~: \texttt{Dim\_Site}, \texttt{Dim\_Temps}, \texttt{Dim\_Motif}, \texttt{Dim\_Individu}, \texttt{Dim\_Mode}, \texttt{Dim\_Menage}, \texttt{Dim\_Geographie}.
\end{itemize}

\begin{figure}[H]
\centering
\includegraphics[width=\textwidth]{modele_constellation.png}
\caption{Schéma en constellation du modèle décisionnel (3 faits, 7 dimensions)}
\label{fig:constellation}
\end{figure}

\section{Description des tables de faits}

\subsection{Fait\_Comptage}
Cette table enregistre les mesures issues des dispositifs de comptage automatique de trafic, à la granularité du site et de l'heure.
\begin{itemize}
\item \textbf{Clés étrangères}~: FK\_id\_Site, FK\_id\_temps
\item \textbf{Mesures}~: nb\_vehicules, volume horaire (matin / soir), flux entrant, flux sortant, désagrégation directionnelle, volume par sens
\end{itemize}

\subsection{Fait\_IndividuMenage}
Cette table croise les informations individuelles et celles du ménage, à la granularité de l'individu enquêté.
\begin{itemize}
\item \textbf{Clés étrangères}~: FK\_id\_individu, FK\_id\_menage, FK\_id\_zone\_residence
\item \textbf{Mesures}~: nombre de déplacements journaliers, taux d'équipement en voiture, nombre de voitures, taux de motorisation, mobilité des seniors, mobilité par zone
\end{itemize}

\subsection{Fait\_Deplacement}
Cette table décrit chaque déplacement individuel déclaré, à la granularité du déplacement.
\begin{itemize}
\item \textbf{Clés étrangères}~: FK\_id\_Dep, FK\_id\_motif, FK\_id\_mode, FK\_id\_zone\_residence
\item \textbf{Mesures}~: durée du trajet, heure de départ, heure d'arrivée, coût total, temps de parcours par motif, corrélation temps / coût
\end{itemize}

\section{Description des dimensions}

\begin{table}[H]
\centering
\caption{Dictionnaire des dimensions du modèle décisionnel}
\begin{tabular}{p{2.6cm}p{6.8cm}p{3.6cm}}
\toprule
\textbf{Dimension} & \textbf{Attributs principaux} & \textbf{Granularité} \\
\midrule
Dim\_Site & id\_Site, code\_site, Ouvl\_Soc, catégorie\_unité, DESC\_Site, Sens & Site de comptage \\
Dim\_Temps & id\_temps, année, MOIS, jour, heure, heure début/fin déplacement & Heure / jour \\
Dim\_Motif & id\_motif, code\_motif, motif\_origine & Motif de déplacement \\
Dim\_Individu & id\_individu, num\_individu, sexe, statut matrimonial, profession, possession permis & Individu \\
Dim\_Mode & id\_mode, code\_mode, libellé, catégorie\_mode & Mode de transport \\
Dim\_Menage & id\_menage, num\_menage, statut occupation, type logement, a\_eau, a\_electricité, a\_internet & Ménage \\
Dim\_Geographie & id\_zone\_residence, code\_zone, quartier, strate, département, difficultés & Zone géographique \\
\bottomrule
\end{tabular}
\end{table}

\section{Indicateurs calculés}
Au-delà des mesures directement issues des sources, plusieurs indicateurs calculés ont été intégrés au modèle ou dérivés lors du chargement~: le \textit{score d'accessibilité} (0 à 100), la \textit{vulnérabilité climatique} (0 à 5), l'\textit{impact sur l'éducation}, le \textit{taux de motorisation des ménages} et la \textit{mobilité des seniors}. Ces deux derniers indicateurs n'ont pu être renseignés que partiellement, la donnée source correspondante étant absente ou incomplète dans l'EMD exploitée.

\section{Conclusion}
Le modèle décisionnel en constellation conçu dans ce chapitre constitue le socle structurel de l'entrepôt de données. Il répond directement aux besoins analytiques identifiés au chapitre précédent tout en assurant la cohérence entre les trois processus métier modélisés. Le chapitre suivant détaille la mise en œuvre du processus ETL permettant d'alimenter ce modèle à partir des sources de données brutes.

\clearpage
